% Generated from project outputs. Do not edit by hand.
\begin{table}[!htbp]
\centering
\begin{threeparttable}
\caption{Including nonlabor by subsample: 2021-2024}
\label{tab:descriptive-2021-2024-subsamples-y3}
\small
\setlength{\tabcolsep}{4pt}
\renewcommand{\arraystretch}{1.15}
\begin{tabularx}{\linewidth}{@{}Xrrrrr@{}}
\toprule
Sample & N & Mean t & SD t & Mean t+1 & SD t+1 \\
\midrule
General & 276 & 3,004.7 & 1,607.3 & 3,355.5 & 1,685.8 \\
Men & 276 & 3,391.6 & 1,781.3 & 3,820.6 & 1,828.7 \\
Women & 276 & 2,390.3 & 1,609.4 & 2,673.7 & 1,716.7 \\
Urban & 276 & 3,215.1 & 1,666.3 & 3,546.3 & 1,717.5 \\
Rural & 276 & 2,669.5 & 1,734.2 & 3,013.7 & 1,930.7 \\
Indigenous & 276 & 2,654.1 & 1,875.4 & 3,132.6 & 1,944.0 \\
Non-Indigenous & 184 & 3,147.9 & 1,648.9 & 3,392.8 & 1,631.1 \\
Bottom 99 percent & 276 & 2,811.3 & 1,363.8 & 3,145.4 & 1,408.9 \\
Bottom 90 percent & 276 & 2,324.4 & 934.6 & 2,590.0 & 936.9 \\
Bottom 50 percent & 276 & 1,338.9 & 537.8 & 1,484.7 & 512.3 \\
\bottomrule
\end{tabularx}
\begin{tablenotes}[flushleft]
\footnotesize
\item Monthly income in quetzales. Unweighted descriptive statistics of survey-weighted cohort means. N counts cohort-transition rows, not individuals or independent cohorts. SD is dispersion across cohort means, not the standard error of a mean. Both incomes must be observed. Subsamples overlap; bottom-income groups are cumulative.
\end{tablenotes}
\end{threeparttable}
\end{table}
